% Imporditud: tools/import_eksperimendid.py
% Allikas: Eesti füüsikaolümpiaadi eksperimendi ülesannete
%   kogu (Rael Kalda, 2023), ülesanne Ü118, lahendus L118.
\setAuthor{Erkki Tempel}
\setRound{piirkonnavoor}
\setYear{2019}
\setNumber{P E2}
\setDifficulty{3}
\setTopic{vedelike mehaanika}
\setVahendid{20 ml süstal, tops veega, kuivatuspaber}
\setPunktid{12}
\setAllikas{Eksperimendi kogu Ü118, lahendus lk 93}
\setLahendusseis{toores_ilma_skeemita}

\prob{Süstal}
Leidke tühja süstla mass.

\hint

\solu
Kasutame süstalt laua serval kangina. Täpsema lugemi võtmiseks tuleb süstla skaala asetada allapoole, siis saab laua servalt täpsemini lugemi võtta. Sätime tühja süstla korral süstla kolvi 5 ml peale ning leiame laua serval tühja süstla masskeskme. Edaspidistes arvutustes võime eeldada, et kogu sütla mass asub masskeskmes. Võtame süstlasse Vv = 5 ml vett ning leiame laua serval tasakaalupunkti. Ühel kangi poolel on Vv = 5 ml vett massiga mv = $\rho$Vv = 5 g ning teisel poolel süstla mass ms. Leiame vee keskpunkti kauguse lv toetuspunktist ning süstla masskeskme kauguse ls toetuspunktist. Kangi tasakaalust saame kirja panna seose, millest

leiame süstla massi

mvlv = msls $\Rightarrow$ ms = mv lv =$\approx$ 10 g ls

\textit{Žürii hindamisskeem on kogumikus lk 93; siia ei ole seda ümber kirjutatud.}
\probend
